\documentclass[11pt]{article}
\usepackage[margin=1in]{geometry}
\usepackage[T1]{fontenc}
\usepackage{lmodern}
\usepackage{amsmath,amssymb,amsthm,bm}
\usepackage{booktabs,graphicx,caption}
\usepackage[authoryear,round]{natbib}
\usepackage[hidelinks]{hyperref}
\usepackage{xcolor}
\usepackage{setspace}
\onehalfspacing

\newtheorem{proposition}{Proposition}
\newtheorem{definition}{Definition}
\theoremstyle{remark}
\newtheorem{remark}{Remark}

\newcommand{\1}{\mathbf{1}}
\newcommand{\E}{\mathbb{E}}
\newcommand{\F}{\mathcal{F}}
\newcommand{\VaR}{\mathrm{VaR}}
\newcommand{\ES}{\mathrm{ES}}
\newcommand{\todo}[1]{\textcolor{red}{[TODO: #1]}}

\title{\todo{Title after the holdout evaluation}}
\author{Mayur Agrawal\thanks{Department of Computer Science and Engineering and Department of Economic Sciences, Indian Institute of Technology Kanpur, Kanpur 208016, India. Email: \href{mailto:mayurag22@iitk.ac.in}{mayurag22@iitk.ac.in}. Code, data pipeline and the pre-registered protocol: \url{https://github.com/mayur-agrawal-iitk/learning-the-tail}.}}
\date{\today}

\begin{document}
\maketitle

\begin{abstract}
Does deep learning improve forecasts of Value-at-Risk and Expected Shortfall when it is given more information, or only when it is given a better architecture? We study pre-open forecasts for 24 equity indices, using the returns of foreign markets that close before each market opens. All models are trained or evaluated with the strictly consistent FZ0 score. A pre-registered protocol develops the models on twelve indices and tests the frozen models exactly once on twelve further indices. Cross-market information is valuable: adding two overnight-news terms to a GJR-GARCH model lowers the loss at every risk level (Diebold--Mariano $t\approx-3.9$ in both samples), with the largest gains in Asian markets after U.S. and European sell-offs. Neural networks, however, do not extract this information. An attention-based contagion network, a neural news-impact GARCH with a proven stability property, and versions trained with an efficiency-improving composite score all fail to beat the two-parameter linear model. An ablation shows that the contagion network gains essentially nothing from the news, which the linear model turns into a large improvement. Forecast combinations that include both families are the most accurate in both samples and spend the fewest days in the Basel yellow zone. In tail-risk forecasting from daily data, information and inductive bias matter more than architecture.

\end{abstract}

\noindent\textbf{Keywords:} Expected Shortfall, Value-at-Risk, volatility spillovers, contagion, deep learning, elicitability, pre-registration.\\
\noindent\textbf{JEL:} C45, C52, C53, C58, G15, G17.

%==============================================================================
\section{Introduction}\label{sec:intro}
%==============================================================================
\todo{intro: written after the holdout evaluation.}


%==============================================================================
\section{Information set: what is known before a market opens}\label{sec:info}
%==============================================================================
Equity markets in different time zones trade in sequence. When the risk manager of market $i$ sets the Value-at-Risk for trading day $D$ shortly before $i$ opens, markets that have traded in the meantime have already revealed information that $i$'s own close-to-close history does not contain \citep{hamao1990correlations,lin1994do}. Let $C_i(D)$ and $O_i(D)$ denote the close and open of market $i$ on its local trading day $D$, in coordinated universal time, and let $D^-$ be $i$'s previous trading day.

\begin{definition}[Pre-open information set]
$\F^{\mathrm{pre}}_{i,D}$ contains the returns of market $i$ up to $D^-$ and, for every other market $j$, all returns of sessions $s$ with $C_j(s)\le O_i(D)-2m$, where $m$ is a safety margin.
\end{definition}

The new information for day $D$ is the overnight move of each foreign market,
\begin{equation}\label{eq:news}
x_{j,D}=\sum_{s:\,C_i(D^-)<C_j(s)\le O_i(D)-2m}\frac{r_{j,s}}{s_{j,s}},
\end{equation}
where $s_{j,s}$ is $j$'s trailing 250-day root mean square return up to the previous session, so that $x_{j,D}$ is measured in units of $j$'s usual volatility. An indicator $a_{j,D}$ records whether at least one such session exists.

Historical exchange schedules have changed (Tokyo extended its close by 30 minutes in 2024, Seoul in 2016, Brazil abolished daylight saving in 2019). To rule out look-ahead bias, we use the \emph{latest} plausible close of every exchange, including closing auctions, and the \emph{earliest} plausible open, including pre-opening sessions. We then require an additional margin of $m=30$ minutes on each side, and handle daylight saving with the IANA time-zone database. A violation could only arise from a schedule error of more than an hour in the conservative direction. The construction makes stale information look new, which can only weaken the predictive content of the news, never strengthen it. In our development sample, Asian markets see the previous U.S.\ and European sessions on about 97\% of days, and U.S.\ markets see the same day's Asian sessions on about 95\% of days.

%==============================================================================
\section{Scoring and statistical efficiency}\label{sec:theory}
%==============================================================================
As in \citet{agrawal2026learning}, forecasts of $(\VaR_\alpha,\ES_\alpha)$ for $\alpha\in\{1\%,2.5\%,5\%\}$ are trained and evaluated with the strictly consistent FZ0 loss of \citet{patton2019dynamic},
\[
L_\alpha(y,v,e)=-\frac{1}{\alpha e}\1\{y\le v\}(v-y)+\frac{v}{e}+\log(-e)-1 .
\]
FZ0 receives information about its target mainly from tail days. Away from the tail, the loss is $v/e+\log(-e)-1$, which depends on the forecast only through its level. Learning \emph{dynamics}, i.e.\ how risk responds to news, from FZ0 alone is therefore statistically inefficient. A Gaussian quasi-likelihood instead learns the dynamics of volatility from every day's squared return. This motivates an auxiliary score.

\begin{proposition}[Composite score]\label{prop:composite}
Let $Q(y,\varsigma^2)=y^2/\varsigma^2+\log\varsigma^2$ (QLIKE). For $\lambda\ge0$, the score
\[
S(y;\{v_k,e_k\}_{k=1}^K,\varsigma^2)=\sum_{k=1}^K L_{\alpha_k}(y,v_k,e_k)+\lambda\,Q(y,\varsigma^2)
\]
is strictly consistent for the vector $\bigl(\{\VaR_{\alpha_k},\ES_{\alpha_k}\}_{k=1}^K,\ \E[y^2]\bigr)$, provided $\E y^2<\infty$ and every $L_{\alpha_k}$ is strictly consistent for its pair (Proposition 1 of \citealp{agrawal2026learning}).
\end{proposition}
\begin{proof}
$\E Q(Y,\varsigma^2)=\E[Y^2]/\varsigma^2+\log\varsigma^2$ has derivative $(\varsigma^2-\E Y^2)/\varsigma^4$ and is uniquely minimised at $\varsigma^2=\E Y^2$. The expected composite score is a sum of terms that depend on disjoint components of the report, each uniquely minimised at the corresponding true functional, so the sum is uniquely minimised at the true vector.
\end{proof}

\begin{remark}
In a location-scale network $v_k=\varsigma q_k$, $e_k=\varsigma c_k$, the reports are linked. The composite score then still identifies the truth when the model is correctly specified. Under misspecification, $\lambda$ trades off tail fit against volatility fit. We fix $\lambda=1$ a priori and do not tune it. QLIKE is also the loss that is robust to noise in the volatility proxy \citep{patton2011volatility}.
\end{remark}

%==============================================================================
\section{Models}\label{sec:models}
%==============================================================================
All models produce one-step-ahead $(\VaR,\ES)$ forecasts at the three levels. The benchmarks of \citet{agrawal2026learning} (historical simulation, GARCH-type models with normal, Student-$t$ and skewed-$t$ innovations, filtered historical simulation, GARCH-EVT, CAViaR-ES and two FZ-GAS models) and its deep models (LSTM-FZ, TRF-FZ, GJR-LSTM) use only the close-to-close information set and are carried over unchanged. Four models are new.

\paragraph{GJRX-skt (econometric, pre-open information).} A GJR-GARCH(1,1,1) with Hansen skewed-$t$ innovations whose variance responds linearly to squared foreign news,
\begin{equation}\label{eq:gjrx}
\sigma_t^2=\omega+(a+g\,\1\{\varepsilon_{t-1}<0\})\varepsilon_{t-1}^2+b\,\sigma_{t-1}^2+d_1\bar X_t+d_2\bar X_t^-,
\end{equation}
where $\bar X_t$ is the average of $x_{j,t}^2$ over markets with fresh news and $\bar X_t^-$ the same average restricted to negative moves. Both are known before $r_t$ is realised. The model is estimated by maximum likelihood with $d_1,d_2\ge0$ on the same rolling window and refit schedule as the other econometric models. It nests GJR-skt ($d_1=d_2=0$).

\paragraph{NG-FZ (neural news-impact GARCH).} Returns are scaled by a per-asset constant, $y_t=r_t/s$, and a $d$-dimensional state follows a \emph{linear} recursion driven by a \emph{learned} news-impact function \citep{engle1993measuring}:
\begin{align}
h_t&=A\,h_{t-1}+N_\theta(y_{t-1}),\qquad A=\mathrm{diag}(a_1,\dots,a_d),\quad 0<a_k<\rho<1,\label{eq:ng}\\
N_\theta(y)&=B\,(y,\,y^2,\,y^2\1\{y<0\},\,|y|)'+c+g\cdot\mathrm{MLP}(\tanh(y/2),\log(1+y^2),\1\{y<0\}).\nonumber
\end{align}
The first coordinate is a GJR-GARCH variance, the others are volatility memories at decay rates $0.5,\dots,0.995$ (as in component GARCH). Log-variance and tail shape are read out by small networks, and the output uses the monotone head of \citet{agrawal2026learning}.

\begin{proposition}[Stability and nesting]\label{prop:ng}
(i) For any input path and any two initial states, $\lVert h_t-h'_t\rVert_\infty\le\rho^{\,t}\lVert h_0-h'_0\rVert_\infty$: the filter forgets its initialisation exponentially (invertibility).
(ii) If $(y_t)$ is strictly stationary and ergodic with $\E\log^+\lVert N_\theta(y_0)\rVert<\infty$, then $h^*_t=\sum_{j\ge0}A^jN_\theta(y_{t-1-j})$ converges almost surely, is the unique stationary ergodic solution of \eqref{eq:ng}, and $\lVert h_t-h^*_t\rVert\to0$ exponentially almost surely \citep{brandt1986stochastic,bougerol1993kalman}.
(iii) With $g=g'=0$ (both neural gates closed) and $d=1$, NG-FZ is exactly GJR-GARCH(1,1) with a constant standardised tail.
\end{proposition}
\begin{proof}
(i) $h_t-h'_t=A^t(h_0-h'_0)$ because the input enters additively and does not depend on the state; $\lVert A^t\rVert_\infty=\max_k a_k^t\le\rho^t$. (ii) By the root test, $\limsup_j\lVert A^jN_\theta(y_{t-1-j})\rVert^{1/j}\le\rho\cdot\limsup_j\exp(\log^+\lVert N_\theta(y_{t-1-j})\rVert/j)=\rho<1$ almost surely, because $\E\log^+\lVert N_\theta(y_0)\rVert<\infty$ implies $\log^+\lVert N_\theta(y_{-j})\rVert/j\to0$ (Borel--Cantelli). Uniqueness and the convergence of $h_t$ follow from (i). (iii) By substitution.
\end{proof}
The network is initialised at (iii), with the GARCH part at a fitted GJR model and the tail at the empirical tail of its residuals, so that training starts from a strong econometric model and learns corrections by minimising FZ0. Gradient descent on a linear recursion with bounded decays cannot produce explosive filters, which is a practical advantage over generic recurrent networks.

\paragraph{CTG-FZ (contagion network).} CTG-FZ adds a cross-market module to LSTM-FZ, with identical inputs and hyperparameters otherwise. Each foreign market contributes a token $k_j=\mathrm{MLP}(x_j,|x_j|,\mathrm{emb}_j)$. A query is formed from the target market's embedding and its LSTM state, and attention over the markets with fresh news yields a contagion vector $c=\sum_jw_jk_j$, which is zero when no foreign market traded overnight. Because the tokens are non-linear in $x_j$ and the query depends on the target's own volatility state, the module can represent disproportionate responses to large foreign crashes, market-specific lead--lag structure and state-dependent contagion. The weights $w_j$ form an interpretable transmission matrix. In the final specification (v2, Section~\ref{sec:design}) the output is location--scale, $\VaR=s\,\varsigma\,q$ and $\ES=s\,\varsigma\,c$, with $\varsigma$ from a variance head trained by the composite score of Proposition~\ref{prop:composite}.

\paragraph{LSTMQ-FZ (ablation).} LSTMQ-FZ is identical to CTG-FZ v2 with the foreign news switched off. The comparison CTG-FZ versus LSTMQ-FZ therefore isolates the value of the cross-market information for the network, holding architecture and objective fixed.

\paragraph{Combinations.} COMB-All (all econometric models except HS, plus the three deep models of \citealp{agrawal2026learning}) is extended to COMB-All+X by adding GJRX-skt and CTG-FZ, and to COMB-All+NG by adding NG-FZ. All combinations are equal-weight averages, specified before any result was seen.


%==============================================================================
\section{Pre-registered design}\label{sec:design}
%==============================================================================
\paragraph{Samples.} The \emph{development} sample consists of the twelve indices of \citet{agrawal2026learning}: S\&P 500, Nasdaq Composite, Russell 2000, FTSE 100, DAX, CAC 40, Nikkei 225, Hang Seng, BSE Sensex, S\&P/ASX 200, KOSPI and Bovespa. The \emph{holdout} sample consists of twelve further indices: S\&P/TSX, SMI, TAIEX, Shanghai Composite, IPC, AEX, IBEX 35, KLCI, Jakarta Composite, Straits Times, ATX and BEL 20. The holdout markets were chosen by a mechanical rule, before any model was run on them. The rule requires at least 2{,}000 daily observations before 2006, no multi-month gaps and fewer than fifteen stale prices; it excluded OMXS30, Merval and FTSE MIB. Both samples use the out-of-sample period January 2006 to September 2026, a 2{,}000-day rolling window and the evaluation methods of \citet{agrawal2026learning}.

\paragraph{Protocol.} The protocol, its amendments and every development result are recorded in a public, time-stamped file (\texttt{docs/PROTOCOL.md} in the repository). New models were designed and revised on the development sample only. Before any holdout run, the final specifications were frozen in a git commit (\texttt{7a58b26}, 15 October 2026), and the holdout sample was then evaluated exactly once. The pre-registered hypotheses, at $\alpha=2.5\%$ with panel Diebold--Mariano tests, were:
\begin{itemize}
\item[H1] CTG-FZ has lower FZ0 loss than GJRX-skt (learning contagion versus a linear spillover model with the same information);
\item[H2] CTG-FZ has lower loss than LSTM-FZ (value of the cross-market information for the network);
\item[H3] COMB-All+X has lower loss than COMB-All.
\end{itemize}

\paragraph{Development history.} For transparency we report every variant tried on the development sample (Table~\ref{tab:devlog}). NG-FZ v1 shared one GARCH block across markets. v2 gave each market its own block and overfitted. The diagnosis, that FZ0 learns volatility dynamics inefficiently, motivated v3 (v2 plus the auxiliary QLIKE score of Proposition~\ref{prop:composite}), which roughly halved v2's deficit. Under the pre-specified rule (lowest development loss), v1 was frozen. CTG-FZ v1 reacted to foreign news but did not improve on LSTM-FZ. By the same efficiency diagnosis, v2 adopted the location--scale output with the composite score, and the ablation LSTMQ-FZ was added at that point. A bug in the first v2 run (a variable that silently disabled the new settings) was detected before any result was seen, fixed with a test, and the run was repeated. v2 was frozen by the selection rule.

\begin{table}[htbp]\centering\footnotesize
\caption{Development history (development sample, $\alpha=2.5\%$): FZ0 loss relative to the stated comparison}\label{tab:devlog}
\begin{tabular}{llrr}\toprule
Variant & Comparison & Loss difference & DM $t$ \\\midrule
NG-FZ v1 (pooled GARCH block) & GJR-skt & $+0.0044$ & 0.42 \\
NG-FZ v2 (market-specific block) & GJR-skt & $+0.0272$ & 2.18 \\
NG-FZ v3 (v2 + QLIKE) & GJR-skt & $+0.0146$ & 1.65 \\
CTG-FZ v1 & GJRX-skt / LSTM-FZ & $+0.0177$ / $+0.0016$ & 1.30 / 0.96 \\
CTG-FZ v2 & GJRX-skt / LSTM-FZ & $+0.0153$ / $-0.0007$ & 1.22 / $-0.29$ \\
CTG-FZ v2 & LSTMQ-FZ (ablation) & $+0.0002$ & 0.48 \\
GJRX-skt & GJR-skt & $-0.0193$ & $-3.90$ \\
\bottomrule\end{tabular}
\par\smallskip\parbox{\linewidth}{\scriptsize Frozen for the holdout: NG-FZ v1, CTG-FZ v2, LSTMQ-FZ, GJRX-skt and all models and combinations of \citet{agrawal2026learning}.}
\end{table}


%==============================================================================
\section{Results}\label{sec:results}
%==============================================================================
\todo{results: written after the holdout evaluation.}


%==============================================================================
\section{Conclusion}\label{sec:conclusion}
%==============================================================================
Giving tail-risk models the information of markets that have already closed improves forecasts substantially, and a two-parameter extension of GJR-GARCH captures most of the gain. The deep models we designed for the task did not. These were an attention-based contagion network, a neural GARCH with provable stability, and versions trained with an efficiency-improving composite score. They react to the news, but they do not turn it into better forecasts. The result holds in a development sample and, after a pre-registered freeze, in twelve markets that no model had seen. Combinations of econometric and neural forecasts remain the most accurate and the best calibrated, in line with our companion paper.

The mechanism, as far as our evidence identifies it, is statistical efficiency. FZ0 learns from tail days, while a likelihood learns volatility dynamics from every day. The correct functional form, foreign squared news entering the variance additively, is easy to impose and hard to discover from eight-year windows. The broader lesson for machine learning in finance is that, in low signal-to-noise problems, the inductive bias and the information set deserve as much attention as the architecture. A deep model should be compared with the best structured model \emph{given the same information}, not with a model that has less.

\paragraph{Limitations and extensions.} We use daily close-to-close returns of equity indices and a one-day horizon. Intraday data would sharpen the timing of transmission, and other asset classes would test its generality. Larger pooled samples, for example thousands of individual stocks, may change the efficiency trade-off in favour of flexible models. The finding that tail risk flows mainly from the West to Asia overnight invites a structural follow-up: whether the transmitted component of tail risk is priced, in the spirit of \citet{kelly2014tail} and \citet{bali2009intertemporal}.


\section*{Acknowledgments and disclosure}
The author used AI-based assistance (large language model tools) for code development, literature search and editing. All modelling choices, results and interpretations were checked by the author, who takes full responsibility for the content of this paper. Any errors are the author's own.

\bibliographystyle{plainnat}
\bibliography{references}

\end{document}
